%%%%%%%%%%%%%%%%%%%%%%
%%%%%%%%%%%%%%%%%%%%%%

\begin{figure*}
\begin{centering}

	\begin{subfigure}[b]{0.40\textwidth}
                \centering               
                \includegraphics[width=\textwidth]{Figures/oc/new_results/human_mouse_mfo_4}
                \caption{Human/Mouse}                 
                \label{fig:ocHMM4}
        \end{subfigure}%
	\quad
	\begin{subfigure}[b]{0.40\textwidth}
                \centering               
                \includegraphics[width=\textwidth]{Figures/oc/new_results/human_rat_mfo_4}
                \caption{Human/Rat}                 
                \label{fig:ocHRM4}
        \end{subfigure}%
        
	\begin{subfigure}[b]{0.40\textwidth}
                \centering              
                \includegraphics[width=\textwidth]{Figures/oc/new_results/mouse_rat_mfo_4}
                \caption{Mouse/Rat}                  
                \label{fig:ocMRM4}
        \end{subfigure}%
       \quad
	\begin{subfigure}[b]{0.40\textwidth}
                \centering            
                \includegraphics[width=\textwidth]{Figures/oc/new_results/worm_fly_mfo_4}
                \caption{Worm/Fly}                    
                \label{fig:ocWFM4}
        \end{subfigure}%
       


\par\end{centering}
\caption[Comparison of orthologs and different subtypes of paralogs in four pairs of organisms for Molecular Function annotations]{Comparison of orthologs and different subtypes of paralogs in four pairs of organisms for Molecular Function annotations.  Results are divided into orthologs (red), inparalogs (green), within-species  outparalogs (blue), between-species outparalogs (purple) for four pairs of organisms. Standard error bars are shown.}
\label{fig:oc2_mfo}
\end{figure*}
%%%%%%%%%%%%%%%%%%%%%%
%%%%%%%%%%%%%%%%%%%%%%
